\exercisegroup

\begin{enumerate}
\def\labelenumi{\arabic{enumi})}
\item
  Let I be an uncountable set, and let E be the space \(\mathbf{R}^{(I)}\) endowed with the topology defined in I, p.~24, exerc. 14; let \(E'\) be its dual (IV, p.~50, exerc. 16, c)). Show that in \(E'\) there exist non relatively compact (for \(\sigma(E', E)\) ) subsets H such that from every sequence of points of H we can extract a sequence which converges to a point of H for \(\sigma(E', E)\) .
\item
  \begin{enumerate}
  \def\labelenumii{\alph{enumii})}
  \tightlist
  \item
    Let X be a regular space and A a subset of X. Suppose that every sequence of points of A has a limit point in X and that there exists a metrizable topology T on X which is coarser than the given topology \(T_{0}\) . Show that the closure \(\overline{A}\) of A in X is a compact metrizable space (arguing by reductio ad absurdum, show that the topologies induced by T and \(T_{0}\) on \(\overline{A}\) are identical).
  \end{enumerate}
\end{enumerate}

\begin{enumerate}
\def\labelenumi{\alph{enumi})}
\setcounter{enumi}{1}
\item
  Let E be a Hausdorff locally convex space, which is the union of a sequence \((\mathrm{E}_{n})\) of metrizable vector subspaces for the topology induced by that of E. Show that, for a subset A of E to be relatively compact in E, it is necessary and sufficient that every sequence of points of A has a limit point in E. (If \((\mathrm{W}_{mn})\) (for \(m \geqslant 1\) ) is a fundamental system of convex neighbourhoods of 0 in \(E_{n}\) which are open in \(E_{n}\) , let \(U_{mn}\) be a convex neighbourhood of 0 in E such that \(E_{n} \cap U_{mn} = W_{mn}\) (II, p.~33, lemma 2); on E consider the topology for which the \(U_{mn}\) ( \(m \geqslant 1\) , \(n \geqslant 1\) ) form a fundamental system of neighbourhoods of 0).
\item
  Extend Šmulian's theorem to a strict inductive limit space (II, p.~33) of Fréchet spaces.
\end{enumerate}

\begin{enumerate}
\def\labelenumi{\arabic{enumi})}
\setcounter{enumi}{2}
\tightlist
\item
  \begin{enumerate}
  \def\labelenumii{\alph{enumii})}
  \tightlist
  \item
    Let E be a Hausdorff locally convex space and \(E'\) its dual; suppose that there exists a countable everywhere dense set in \(E'\) for the topology \(\sigma(E', E)\) . Show that the topology of E (resp. \(\sigma(E, E')\) ) is finer than a metrizable locally convex topology.
  \end{enumerate}
\end{enumerate}

\begin{enumerate}
\def\labelenumi{\alph{enumi})}
\setcounter{enumi}{1}
\item
  Let \((x_{n})\) be a sequence of points of E such that every sequence extracted from \((x_{n})\) has a limit point for the initial topology (resp. the topology \(\sigma (\mathrm{E},\mathrm{E}^{\prime})\) ). Show that there exists a sequence extracted from \((x_{n})\) which converges in E for the initial topology (resp. the topology \(\sigma (\mathrm{E},\mathrm{E}^{\prime})\) ). (Let \((a_n^{\prime})\) be an everywhere dense sequence in \(\mathbf{E}'\) for \(\sigma (\mathrm{E}',\mathrm{E})\) ; extract a sequence \((y_{n})\) from \((x_{n})\) such that \((\langle y_n,a_p'\rangle)\) tends to a limit for every index \(p\) , and show that the sequence \((y_{n})\) has only one limit point for the initial topology (resp. for \(\sigma (\mathrm{E},\mathrm{E}^{\prime}))\) .)
\item
  For a subset A of E to be relatively compact for the initial topology (resp. for \(\sigma(E, E')\) , it is necessary and sufficient that from every sequence \((x_n)\) of points of A, we can extract a sequence \((x_{n_k})\) which converges to a point of E for the initial topology (resp. for \(\sigma(E, E')\) ) (use exerc. 2, a)).
\end{enumerate}

\begin{enumerate}
\def\labelenumi{\arabic{enumi})}
\setcounter{enumi}{3}
\item
  Let \(\mathrm{E}\) be the Banach space \(l^{\infty}(\mathbb{N})\), which does not satisfy the first axiom of countability (I, p.~25, exerc. 1), and let \(\mathrm{E}'\) be its dual. For every integer \(n \geqslant 0\) , let \(e_n'\) be the continuous linear form on E which associates to each \(x = (\xi_n) \in \mathrm{E}\) the \(n\) th term of this sequence. Show that the sequence \((e_n')\) is total in \(\mathrm{E}'\) for \(\sigma(\mathrm{E}', \mathrm{E})\) ; moreover, every sequence extracted from \((e_n')\) has a limit point in \(\mathrm{E}'\) , for the topology \(\sigma(\mathrm{E}', \mathrm{E})\) , but there exists no sequence extracted from \((e_n')\) which converges in \(\mathrm{E}'\) for this topology.
\item
  \begin{enumerate}
  \def\labelenumii{\alph{enumii})}
  \tightlist
  \item
    Let X be a compact space, H an arbitrary subset of the space \(\mathcal{C}(\mathrm{X})\) of continuous numerical functions on X. Let \(f_{0}\) be a point of \(\mathcal{C}(\mathrm{X})\) which is in the closure of H for the topology \(T_{s}\) of simple convergence on \(\mathcal{C}(\mathrm{X})\) . Show that there exists a countable subset \(H_{0}\) of H such that \(f_{0}\) is in the closure of \(H_{0}\) for \(T_{s}\) . (Show that for every pair of integers m \textgreater{} 0, n \textgreater{} 0 there exists a finite subset \(\mathrm{H}(m, n)\) of H with the following property: for every set of m points \(t_{k}\) in X ( \(1 \leqslant k \leqslant m\) ), there exists \(f \in \mathrm{H}(m, n)\) such that \(|f_{0}(t_{k}) - f(t_{k})| \leqslant 1/n\) for \(1 \leqslant k \leqslant m\) . b) Let E be a locally convex metrizable space, \(E'\) its dual and H a subset of E. Show that if \(x_{0}\) is in the closure of H for the weakened topology \(\sigma(\mathrm{E}, \mathrm{E}')\) , then there exists a countable subset \(H_{0}\) of H such that \(x_{0}\) is in the closure of \(H_{0}\) for this topology. (Use a), observing that \(E'\) is the union of a countable family of compact sets for \(\sigma(\mathrm{E}', \mathrm{E})\) .)
  \end{enumerate}
\end{enumerate}

¶ 6) a) Let X be a compact space, H a convex subset of the product space \(\mathbf{R}^{\mathrm{X}}\) , consisting of continuous functions on X. Suppose that every decreasing sequence of non-empty convex and closed subsets in H has a non-empty intersection. Show that the closure \(\overline{\mathrm{H}}\) of H in \(\mathbf{R}^{\mathrm{X}}\) is compact and consists of continuous functions on X. (Argue by reductio ad absurdum: consider a non continuous function \(u \in \overline{\mathrm{H}}\) ; show that then there will exist a point \(a \in \mathrm{X}\) , a number \(\delta > 0\) , a sequence \((x_{n})\) points of X and a sequence \((f_{m})\) of functions in H such that:

\(1^{\circ}\left|u(x_{n}) - u(a)\right|\geqslant \delta\) for all \(n;2^{\circ}\left|f_{m}(x_{n}) - f_{m}(a)\right|\leqslant \frac{\delta}{8}\) for \(m\leqslant n;3^{\circ}\left|u(x_{n}) - f_{m}(x_{n})\right|\leqslant \frac{\delta}{8}\) and

\(\left|u(a) - f_m(a)\right| \leqslant \frac{\delta}{8}\) for \(m \geqslant n + 1\) . Consider a limit point \(b\) of the sequence \((x_n)\) , and a function \(f\) belonging to the intersection of the \(\mathbf{A}_m\) , where \(\mathbf{A}_m\) is the closed convex envelope in \(\mathbf{H}\) of the set of all \(f_k\) for \(k \geqslant m\) .

\begin{enumerate}
\def\labelenumi{\alph{enumi})}
\setcounter{enumi}{1}
\tightlist
\item
  Let E be a Hausdorff and quasi-complete locally convex space, \(E'\) its dual. Let H be a convex subset of E such that every decreasing sequence of non-empty and closed convex subsets of H has a non-empty intersection; show that H is relatively compact in E for the topology \(\sigma(\mathrm{E}, \mathrm{E}')\) . (Reduce to the case where E is complete; consider E to be embedded in \(E'^{*}\) and use a) and also III, p.~21, cor. 1.)
\end{enumerate}

\begin{enumerate}
\def\labelenumi{\arabic{enumi})}
\setcounter{enumi}{6}
\item
  Let \(\mathbf{E}\) be a Fréchet space satisfying the first axiom of countability, \(\mathbf{E}'\) its dual. Show that for a convex subset \(\mathbf{A}'\) of \(\mathbf{E}'\) to be closed for \(\sigma(\mathbf{E}', \mathbf{E})\) , it is sufficient that, if a sequence \((x_n')\) of points of \(\mathbf{A}'\) has a limit \(a'\) in \(\mathbf{E}'\) for \(\sigma(\mathbf{E}', \mathbf{E})\) , then \(a' \in \mathbf{A}'\) .
\item
  Let \(\mathrm{F}\) and \(\mathrm{G}\) be two vector spaces in separating duality. Show that the properties \(\alpha)\) and \(\beta)\) of IV, p.~52, exerc. 6, a) are also equivalent to the following:
\end{enumerate}

\(\gamma)\) F, with \(\tau (\mathrm{F},\mathrm{G})\) is quasi-complete, and every bounded sequence of points of F has a limit point for \(\sigma (\mathrm{F},\mathrm{G})\) (cf.~IV, p.~35, th. 1);

δ) F, with \(\tau(\mathrm{F},\mathrm{G})\) is quasi-complete, and every decreasing sequence of non-empty closed convex and bounded sets in F has a non-empty intersection (cf.~exerc. 6, b)).

\begin{enumerate}
\def\labelenumi{\arabic{enumi})}
\setcounter{enumi}{8}
\tightlist
\item
  Let E be a Hausdorff and quasi-complete locally convex space; for E to be semi-reflexive, it is necessary and sufficient that every closed vector subspace of E, in which there exists a countable everywhere dense subset, is semi-reflexive (cf.~IV, p.~35, th. 1).
\end{enumerate}

¶ 10) Let E be a Hausdorff, quasi-complete and non semi-reflexive locally convex space, and let H be a closed hyperplane in E containing the origin. Let \((\mathbf{C}_{n})\) be a decreasing sequence of non-empty, closed convex and bounded sets, contained in H and not containing 0 and whose intersection is empty (exerc. 8). Let x be a point not belonging to H, and let A be the closed

convex and balanced envelope of the union of the sets \(\left(1 - \frac{1}{n}\right)x + C_n\) for \(n > 0\) .

\begin{enumerate}
\def\labelenumi{\alph{enumi})}
\item
  Show that there exists no supporting hyperplane of A which is parallel to H (for \(y \in x + H\) , observe that there exists an integer n such that \(y \notin x + C_{n}\) ).
\item
  Let \(z \in \mathbf{H}\) be such that \(z \notin \mathbf{C}_1\) ; show that the convex envelope of the union of two closed convex and bounded sets \(\mathbf{A}\) and \(\mathbf{B} = x + z + \mathbf{C}_1\) , is not closed (prove that \(x + z\) is in the closure of this envelope, but does not belong to it).
\end{enumerate}

\begin{enumerate}
\def\labelenumi{\arabic{enumi})}
\setcounter{enumi}{10}
\item
  Let A be an infinite set, and let \(E = \overline{\mathcal{H}}(A)\) (IV, p.~47, exerc. 1), \(E' = \ell^{1}(A)\) its dual, \(E'' = \ell^{\infty}(A)\) its bidual (IV, p.~54, exerc. 14). Show that every subset of \(E'\) , which is compact for \(\sigma(E', E'')\) is strongly compact (use Šmulian's theorem and IV, p.~54, exerc. 15).
\item
  Let E be a non reflexive Banach space. Prove that there exists a closed, non reflexive vector subspace M of infinite codimension in E. (Let \((x_{n})\) be a bounded sequence in E which has no limit point for the topology \(\sigma(\mathrm{E},\mathrm{E}')\) (IV, p.~68, exerc. 8); by induction construct a sequence \((x_{n_{k}})\) extracted from \((x_{n})\) and a topologically free sequence \((y_{k})\) such that \(\| x_{n_{k}} - y_{k} \| \leqslant 1/k\) for \(k \geqslant 1\) , and consider the closed vector subspace of E generated by the \(y_{2k}\) .
\end{enumerate}

¶ 13) Let E be a non reflexive Banach space satisfying the first axiom of countability, and let M be a non reflexive closed vector subspace of E of infinite codimension (exerc. 12). Let \((x_{n})\) be an everywhere dense sequence in the unit sphere, C the closed convex balanced envelope of the sequence \((x_{n}/n)\) ; C is strongly compact in E and \(C + M = A\) is a closed convex set (GT, III, § 4, No.~1, cor. 1 to prop. 1). Let S be the unit ball in E, and \(B = A \cap S\) .

\begin{enumerate}
\def\labelenumi{\alph{enumi})}
\item
  Show that 0 is not an interior point of A and deduce that there exists \(x_0 \in E\) such that \(\lambda x_0 \notin A\) for \(\lambda > 0\) .
\item
  Show that there exists no supporting hyperplane of B passing through 0 (observe that such a hyperplane must be a supporting hyperplane of C).
\item
  Let \(\mathrm{U}_0 = \mathrm{M} \cap \mathrm{S}\) , and let \((\mathrm{U}_n)_{n \geqslant 1}\) be a decreasing sequence of closed convex, bounded and non-empty sets, such that \(\mathrm{U}_1 \subset \frac{1}{2}\mathrm{U}_0\) and such that the intersection of the \(\mathrm{U}_n\) is empty (IV, p.~68, exerc. 8). Let \(\mathbf{F}\) be the closed convex envelope of the union of the sets \(\frac{1}{n} x_0 + \mathrm{U}_n\) ( \(n \geqslant 1\) ). Show that \(\mathbf{B} \cap \mathbf{F} = \emptyset\) but that there exists no closed hyperplane separating \(\mathbf{B}\) and \(\mathbf{F}\) (use \(b\) )).
\end{enumerate}

\begin{enumerate}
\def\labelenumi{\arabic{enumi})}
\setcounter{enumi}{13}
\tightlist
\item
  Let \(\mathbf{E}\) be a Banach space satisfying the first axiom of countability. A sequence \((e_n)_{n\geqslant 0}\) of elements of \(\mathbf{E}\) is said to be a Banach basis if the following property is satisfied: for every \(x\in \mathbb{E}\) , there exists a unique sequence \((a_{n})\) of scalars such that \(x = \sum_{n = 0}^{\infty}a_{n}e_{n}\) , where the series on
\end{enumerate}

the right hand side is convergent.

\begin{enumerate}
\def\labelenumi{\alph{enumi})}
\tightlist
\item
  Show that the family \((e_n)\) is total and free. Let \(E_n\) be the vector subspace (closed) of E generated by the \(e_m\) for indices \(m \leqslant n\) , and let \(P_n\) be the projector from E onto \(E_n\) defined by \(P_n \cdot (\sum_{m=0}^{\infty} a_m e_m) = \sum_{m=0}^{n} a_m e_m\) . Show that the \(P_n\) are continuous linear mappings and that
\end{enumerate}

\(\sup_{n}\| P_{n}\| < + \infty .\) (Consider the norm on E defined by \(\left\| \sum_{n = 0}^{\infty}a_n e_n\right\| = \sup_{n}\left\| \sum_{m = 0}^{n}a_me_m\right\|\)

show that E is complete for this norm and deduce that it is equivalent to the given norm on E (cf.~I, p.~17, th. 1).) Show that for every pair of integers \(p < q\) , the norm of the projection \(P_{q,p}\) of \(\mathrm{E}_q\) onto \(\mathrm{E}_p\) , which is parallel to the vector subspace generated by the \(e_m\) for indices \(m\) such that \(p + 1 \leqslant m \leqslant q\) , is bounded by a number independent of \(p, q\) .

\begin{enumerate}
\def\labelenumi{\alph{enumi})}
\setcounter{enumi}{1}
\tightlist
\item
  Conversely, let \((e_{n})\) be a total sequence of elements of E, which is a free family and is such that the norms of the projections \(P_{q,p}\) for \(0 \leqslant p < q\) are bounded by a number M independent of p and q. Show that \((e_{n})\) is a Banach basis of E. (First prove that the sequence \((e_{n})\) is topologically free and define the projectors \(P_{n}\) ; then \(\|P_{n}\| \leqslant M\) for all n.~Next observe that if \(d(x, E_{n})\) is the distance of a point \(x \in E\) from \(E_{n}\) , then \(\|x - P_{n} \cdot x\| \leqslant (M + 1) d(x, E_{n})\) .) \(^{1}\)
\end{enumerate}

\(^{1}\) It is clear that the existence of a Banach basis in E implies that E satisfies the first axiom of countability. But there are examples of Banach spaces satisfying the first axiom of countability in which there does not exist a Banach basis (P. ENFLO, Acta Math., t. CXXX (1973), p.~309-317).

¶ 15) Let E be a Banach space with a Banach basis \((e_n)\) (exerc. 14); then there exists a unique sequence \((e_n')\) in the dual \(E'\) of E such that the expression of every \(x \in E\) in terms of the Banach basis \((e_n)\) can be written as \(x = \sum_{n=0}^{\infty} \langle x, e_n' \rangle e_n\) .

\begin{enumerate}
\def\labelenumi{\alph{enumi})}
\item
  Let \(\mathrm{F}_n\) be the closed vector subspace of \(\mathbf{E}\) generated by the \(e_m\) for \(m \geqslant n\) . For every \(x' \in \mathbf{E}'\) , let \(\| x' \|_n\) denote the norm of the restriction of the linear form \(x'\) to \(\mathrm{F}_n\) . Show that, for \((e_n')\) to be a Banach basis of \(\mathbf{E}'\) , it is necessary and sufficient that for every \(x' \in \mathbf{E}'\) , the sequence \((\| x' \|_n)\) tends to 0. (Consider the transpose \({}^t P_n\) and evaluate the norm \(\| {}^t P_n \cdot x' - x' \|\) .) In this case, the Banach basis \((e_n)\) is said to be contracting.
\item
  Suppose that the Banach basis \((e_n)\) is contracting. Show that for every point \(x''\) of the bidual \(\mathrm{E}''\) of E, the sequence of sums \(\sum_{m=0}^{n} \langle e_m', x'' \rangle e_m\) is bounded in E (consider the transpose \(t^t P_n\) )
\end{enumerate}

in \(\mathrm{E}''\) ). Conversely, for every sequence \((a_{n})\) of scalars such that the sequence of sums \(\sum_{m=0}^{n} a_{m} e_{m}\)

is bounded in E, there exists a unique \(x'' \in E''\) such that \(\langle e_{n}', x'' \rangle = a_{n}\) for all n (use the compactness of a closed ball in \(E''\) for the topology \(\sigma(E'', E')\) ).

\begin{enumerate}
\def\labelenumi{\alph{enumi})}
\setcounter{enumi}{2}
\item
  A Banach basis \((e_n)\) in E is said to be complete if, for every sequence \((a_n)\) of scalars such that the sequence of sums \(\sum_{m=0}^{n} a_m e_m\) is bounded, the series \(\sum_{n=0}^{\infty} a_n e_n\) converges. If \((e_n)\) is a contracting basis of E, the basis \((e_{n}^{\prime})\) of \(E^{\prime}\) is complete (use the compactness of a closed ball in \(E^{\prime}\) for \(\sigma(\mathrm{E}^{\prime},\mathrm{E})\) ). d) In general, the sequence \((e_{n}^{\prime})\) is a Banach basis of the closed subspace \(F^{\prime}\) of the strong dual \(E^{\prime}\) of E, generated by the \(e_{n}^{\prime}\) , and there is an injective continuous linear mapping J from E into the strong dual \(F^{\prime\prime}\) of \(F^{\prime}\) such that \(\langle J.x,z^{\prime}\rangle=\langle x,z^{\prime}\rangle\) for all \(x\in E\) and all \(z^{\prime}\in F^{\prime}\) . Show that there exists a constant K\textgreater0 such that \(\|J.x\|\geqslant K.\|x\|\) , and that if the basis \((e_{n})\) is complete, then J is an isomorphism from E onto the topological vector space \(F^{\prime\prime}\) .
\item
  Show that for \(\mathbf{E}\) to be reflexive, it is necessary and sufficient that the basis \((e_n)\) is contracting and complete (use \(b\) ).
\end{enumerate}

\begin{enumerate}
\def\labelenumi{\arabic{enumi})}
\setcounter{enumi}{15}
\tightlist
\item
  \begin{enumerate}
  \def\labelenumii{\alph{enumii})}
  \tightlist
  \item
    Let \(\mathbf{E}\) be a Banach space. For an infinite sequence \((x_{n})\) of points of \(\mathbf{E}\) , the following properties are equivalent:
  \end{enumerate}
\end{enumerate}

\(\alpha)\) The series with the general term \(x_{n}\) is commutatively convergent (GT, III, § 5, No.~7).

\(\beta)\) For every subset I of \(\mathbf{N}\) , the series defined by the sequence \((x_{n})_{n\in \mathrm{I}}\) is convergent (GT, III, § 5, No.~3, prop. 2 and § 5, exerc. 4).

\(\gamma)\) For every sequence \((\varepsilon_{n})\) of numbers equal to 1 or to \(-1\) , the series with the general term \(\varepsilon_{n}x_{n}\) is convergent.

\(\delta)\) For every \(\varepsilon > 0\) , there exists a finite subset \(J\) of \(N\) such that, for every finite subset \(H\) of \(N\) not intersecting \(J\) , we have \(\| \sum x_n \| \leqslant \varepsilon\) .

\begin{enumerate}
\def\labelenumi{\alph{enumi})}
\setcounter{enumi}{1}
\tightlist
\item
  Let \((e_n)\) be a Banach basis of E. The following properties are equivalent:
\end{enumerate}

\(\alpha)\) For every permutation \(\pi\) of \(\mathbf{N}\) , the sequence \((e_{\pi(n)})\) is a Banach basis of E.

\(\beta)\) For every sequence \((\varepsilon_{n})\) of numbers equal to 1 or to -1, the sequence \((\varepsilon_{n}e_{n})\) is a Banach basis of E.

\(\gamma)\) For every \(x = \sum_{n=0}^{\infty} \xi_n e_n\) in E and every sequence \((\eta_n)_{n \in \mathbb{N}}\) for which \(|\eta_n| \leqslant |\xi_n|\) for all \(n\) , the series with the general term \(\eta_n e_n\) converges in E.

δ) For every \(x = \sum_{n=0}^{\infty} \xi_n e_n\) in E and every strictly increasing sequence \((n_k)_{k \in \mathbf{N}}\) of integers \(\geqslant 0\) , the series with the general term \(\xi_{n_k} e_{n_k}\) converges in E.

ε) There exists a real number M \textgreater{} 0 such that, for every finite subset J of N and for every \(x = \sum_{n=0}^{\infty} \xi_{n} e_{n}\) in E, we have \(\left\|\sum_{n \in J} \xi_{n} e_{n}\right\| \leqslant M \|x\|\) .

(To prove that \(\alpha\) ) implies \(\varepsilon\) ), argue as in IV, p.~69, exerc. 14. To prove that \(\beta\) implies \(\gamma\) ), reduce to the case where the \(\xi_{n}\) and the \(\eta_{n}\) are real, and consider the sums \(\sum_{n=p}^{q} \langle \eta_{n} e_{n}, x' \rangle\) for \(x' \in E'\) .)

When these conditions are satisfied, \((e_n)\) is said to be an unconditional Banach basis of E. c) Suppose \((e_n)\) is an unconditional basis. Show that there exists a real number \(K > 0\) such that, for every sequence \((\varepsilon_n)\) of numbers equal to 1 or -1 and for all \(x = \sum_{n=0}^{\infty} \xi_n e_n\) in E, we have \(\| \sum_{n=0}^{\infty} \varepsilon_n \xi_n e_n \| \leqslant M \| x\|\) (same method as in IV, p.~69, exerc. 14). Deduce that for every

bounded sequence of scalars \((\lambda_{n})\) and for all \(x = \sum_{n=0}^{\infty} \xi_{n} e_{n}\) in X, we have

\[
\left\| \sum_ {n = 0} ^ {\infty} \lambda_ {n} \xi_ {n} e _ {n} \right\| \leqslant 2 K \sup _ {n} | \lambda_ {n} |. \left\| \sum_ {n = 0} ^ {\infty} \xi_ {n} e _ {n} \right\|
\]

(argue as in \(b\) ), for the proof that \(\beta\) implies \(\gamma\) ).

¶ 17) a) Let E be a Banach space with an unconditional Banach basis ( \(e_n\) ) (exerc. 16). Show that if the basis ( \(e_n\) ) is not contracting (IV, p.~70, exerc. 15), then there exists a number \(\alpha > 0\) , a linear form \(x' \in E'\) such that \(\|x'\| = 1\) , a strictly increasing sequence of integers ( \(n_k\) ) and for each \(k\) , an element \(y_k\) which is a linear combination of the \(e_j\) for \(n_k \leqslant j \leqslant n_{k+1}\) , and is such that \(\|y_k\| \leqslant 1\) and \(\langle y_k, x' \rangle \geqslant \alpha\) . Deduce that for every finite sequence ( \(\lambda_j\) ) \(_{1 \leqslant j \leqslant m}\) of scalars, we have \(\|\sum_{j=1}^{m} \lambda_j y_j\| \geqslant \frac{\alpha}{2K} \sum_{j=1}^{m} |\lambda_j|\) (use exerc. 16, c)). Conclude that there exists a topological vector space isomorphism from \(\ell^1 (\mathbf{N})\) onto a closed subspace of E.

\begin{enumerate}
\def\labelenumi{\alph{enumi})}
\setcounter{enumi}{1}
\item
  Deduce from \(a\) ) that if the strong dual \(\mathrm{E}'\) of \(\mathrm{E}\) satisfies the first axiom of countability, then every unconditional basis \((e_n)\) of \(\mathrm{E}\) is contracting, and hence \((e_n')\) is an unconditional basis of \(\mathrm{E}'\) (IV, p.~70, exerc. 15, a)). (Observe that if a closed vector subspace of \(\mathrm{E}\) is isomorphic to \(\ell^1(\mathbf{N})\) , then \(\mathrm{E}'\) cannot satisfy the first axiom of countability.)
\item
  Show that if E has an unconditional basis and if the strong bidual \(E''\) of E satisfies the first axiom of countability, then E is reflexive. (Using IV, p.~51, exerc. 25, note that the strong dual \(E'\) of E satisfies the first axiom of countability; then use IV, p.~70, exerc. 15, c) and IV, p.~53, exerc. 11.)
\end{enumerate}

¶ 18) a) In the space \(\mathbf{R}^{\mathbf{N}}\) of all infinite sequences of real numbers, consider the set J of all sequences \(x = (\xi_n)\) for which the number

\[
\| x \| = \sup \left(\left(\xi_ {p _ {1}} - \xi_ {p _ {2}}\right) ^ {2} + \left(\xi_ {p _ {2}} - \xi_ {p _ {3}}\right) ^ {2} + \dots + \left(\xi_ {p _ {m - 1}} - \xi_ {p _ {m}}\right) ^ {2} + \xi_ {p _ {m + 1}} ^ {2}\right) ^ {1 / 2}
\]

is finite, the supremum being taken for all integers \(m \geqslant 1\) and for all strictly increasing sequences of integers \(p_{1} < p_{2} < \cdots < p_{m+1}\) . Show that \(\|x\|\) is a norm on J and that for this norm J is a Banach space. For every \(x \in E\) , show that the sequence \((\xi_{n})\) has a finite limit \(u(x)\) and that u is a non zero continuous linear form on E; let \(J_{0}\) denote the closed hyperplane of J with equation \(u(x) = 0\) (R. C. James' space).

\begin{enumerate}
\def\labelenumi{\alph{enumi})}
\setcounter{enumi}{1}
\item
  Show that the vectors \(e_n = (\delta_{mn})_{m \geq 0}\) form a Banach basis of \(J_0\) and that this basis is contracting (IV, p.~70, exerc. 15). (To prove the latter point, argue as in exerc. 17, a), by showing that the constructed sequence \((y_k)\) , is such that the series with the general term \(y_k / k\) is convergent in E; this implies a contradiction.)
\item
  Show that the identity mapping from \(\mathrm{J}_0\) onto itself can be extended to a topological vector space isomorphism from the strong bidual \(\mathrm{J}_0''\) onto \(\mathrm{J}\) , in such a way that \(\mathrm{J}_0'' / \mathrm{J}_0\) is 1-dimensional (use IV, p.~70, exerc. 15, b)). Deduce that no Banach basis of \(\mathrm{J}_0\) can be unconditional (exerc. 17, c)).
\end{enumerate}

* \(d\) ) Let \(\mathrm{H}_1,\mathrm{H}_2\) be two closed vector subspaces of \(\mathbf{J}_0\) generated by the \(e_{2n}\) and the \(e_{2n + 1}\) respectively, for \(n\geqslant 0\) . Show that as topological vector spaces, \(\mathrm{H}_{1}\) and \(\mathrm{H}_{2}\) are isomorphic to the Hilbert space \(\ell^2 (\mathbf{N})\) , and that \(\mathbf{J}_0\) is not the sum of \(\mathrm{H}_1\) and \(\mathrm{H}_2.\)

\begin{enumerate}
\def\labelenumi{\alph{enumi})}
\setcounter{enumi}{4}
\tightlist
\item
  Show that on \(\mathrm{J}_0\) there exists no complex locally convex space structure having the real locally convex space structure of \(\mathrm{J}_0\) as the underlying structure (cf.~IV, p.~52, exerc. 3).
\end{enumerate}

